\documentclass[11pt]{article}
% ---------- Anglo-Saxon business letter ----------
\usepackage[T1]{fontenc}
\usepackage[utf8]{inputenc}
\usepackage{lmodern}
\usepackage[english]{babel}
\usepackage[a4paper,top=2cm,bottom=2cm,left=2.4cm,right=2.4cm]{geometry}
\usepackage{microtype}
\usepackage{xcolor}
\usepackage[hidelinks]{hyperref}
\definecolor{accent}{HTML}{1F4E79}
\setlength{\parindent}{0pt}
\setlength{\parskip}{0.75em}
\pagestyle{empty}

% ---------- Sender (edit here) ----------
\newcommand{\myname}{Camille Vedani}
\newcommand{\myaddress}{Toulouse, France}
\newcommand{\myphone}{+33 6 29 27 59 31}
\newcommand{\myemail}{camille.vedani@outlook.fr}

\begin{document}

% ---------- Letterhead ----------
{\LARGE\bfseries\color{accent}\myname}\par
\vspace{-0.5em}{\color{accent}\rule{\linewidth}{0.8pt}}\par
\vspace{-0.6em}{\small \myaddress\enspace$\cdot$\enspace \myphone\enspace$\cdot$\enspace
\href{mailto:camille.vedani@outlook.fr}{\myemail}\\
\href{https://www.linkedin.com/in/camille-vedani-280967270}{linkedin.com/in/camille-vedani-280967270}\enspace$\cdot$\enspace
\href{https://github.com/camved}{github.com/camved}}\par
\vspace{1em}

\today

Malin Vestin, International Coordinator\\
School of Electrical Engineering and Computer Science (EECS)\\
KTH Royal Institute of Technology\\
SE-100 44 Stockholm, Sweden

\textbf{Re: Application for a six-month study period in dependable embedded systems}

Dear Malin Vestin,

% ---------------- ME ----------------
I am a second-year engineering student at ENSEEIHT (Toulouse INP), in the Computer Science \& Networks
apprenticeship track (FISA), and an apprentice software engineer at Kyber, a French open-source start-up building
infrastructure for real-time interactions. At Kyber, I develop an SDK that enables remote control of
machines with less than 20\,ms of latency: I implement network protocols based on QUIC and WebTransport
and optimise our Rust/WebAssembly codebase for browser integration. My path is deliberately broad: two
years of preparatory classes in life sciences (CPGE BCPST), a double bachelor's degree in mathematics and
computer science at Aix-Marseille University, and a first master's year in embedded and connected systems
at Toulouse III. Two research internships, on state reduction for probabilistic automata at LIS (CNRS) and
on execution decision diagrams for worst-case execution time analysis at IRIT, taught me the rigour of
academic research. Projects such as a DPLL SAT solver and model checking with Promela gave me a strong
taste for formal reasoning. All of this serves one goal: designing safe, open medical embedded devices
such as insulin pumps.

% ---------------- YOU ----------------
KTH Royal Institute of Technology stands out for its research and teaching in embedded systems with a clear
focus on dependability and formal methods: exactly the combination I believe medical devices require. Its
close ties with industrial leaders such as Ericsson and Scania, and its strength in medical technology,
bring together the three fields I want to connect: networks, critical embedded software and health. The
Swedish approach to engineering, based on collaboration, autonomy and openness, also strongly matches the
way I already work at Kyber.

% ---------------- US ----------------
I would like to spend one semester (six months) at your university, within the exchange agreement currently being established between ENSEEIHT (Toulouse INP) and KTH. I would bring
my skills in Rust, C/C++, Python and model checking, my experience of real-time networking, and the autonomy I have developed while working fully remotely in an international,
English-speaking team. In return, this semester would let me deepen my knowledge of dependable systems and formal verification in a world-class academic
environment while fulfilling the
international mobility required by my engineering degree. I would be glad to provide any further
information about my academic record. My CV and portfolio are attached.

Thank you for considering my application. I look forward to hearing from you.

Sincerely,

\vspace{1.8em}
\myname
\end{document}
